%========
% WG Edit
% SRI LO: generated spectra and consistency relations at relative order p=0.
%========
\documentclass[10pt]{article}
\usepackage[T1]{fontenc}
\usepackage{lmodern,amsmath,amssymb,geometry}
\geometry{a4paper,margin=2cm}
\allowdisplaybreaks[4]
\setlength{\jot}{4pt}
\begin{document}
\begin{center}\Large SRI (k-inflation): LO\end{center}
The spectra are evaluated at $-k_*\eta_*=1$, through relative amplitude
order $p=0$. The Hubble-flow parameters are denoted by $\epsilon_j$
and $M_{\rm Pl}$ is the reduced Planck mass. Coefficients are generated
numerically from the normalized mode equations, using the iterative algebra.
The convention for the spectral derivatives is
\[
\alpha_n^A=\frac{d^{n+1}\ln P_A}{d(\ln k)^{n+1}},\qquad
\alpha_0^s=n_s-1,\qquad\alpha_0^t=n_t,\qquad r=P_t/P_s.
\]
For k-inflation, $c_t=1$. The scalar amplitude includes transport from the acoustic reference to the common tensor reference. The sound speed is $c_{s*}>0$, with
\[\alpha_n^{c_s}=\frac{d^{n+1}\ln c_s}{d(\ln k)^{n+1}}.\]
\section{Spectra}

\subsection{Scalar spectrum}
\[P_s(k_*)=\frac{H_*^2}{8\pi^2M_{\rm Pl}^2\epsilon_{1*}c_{s*}}\left[\sum_{d=0}^{0}F_{s,d}+O(\mathrm{SR}^{1})\right].\]
\paragraph{Relative order 0.}
\begin{align*}
F_{s,0} &= 1
\end{align*}

\subsection{Tensor spectrum}
\[P_t(k_*)=\frac{2H_*^2}{\pi^2M_{\rm Pl}^2}\left[\sum_{d=0}^{0}F_{t,d}+O(\mathrm{SR}^{1})\right].\]
\paragraph{Relative order 0.}
\begin{align*}
F_{t,0} &= 1
\end{align*}

\section{Consistency relations}
\subsection{$n=0$}
For $p=0$, the absolute cutoff is $W=1$:
\[\alpha_{0}^t=\sum_{d=1}^{1}R_{0,d}+O(\mathrm{SR}^{2}).\]
\paragraph{Absolute order 1.}
\begin{align*}
R_{0,1} &=  - \frac{r}{8 c_{s}}
\end{align*}

\subsection{$n=1$}
For $p=0$, the absolute cutoff is $W=2$:
\[\alpha_{1}^t=\sum_{d=2}^{2}R_{1,d}+O(\mathrm{SR}^{3}).\]
\paragraph{Absolute order 2.}
\begin{align*}
R_{1,2} &= \frac{r^{2}}{64 c_{s}^{2}} + \frac{\alpha_{0}^{c_s} r}{8 c_{s}} \\
&{}  + \frac{\alpha_{0}^s r}{8 c_{s}}
\end{align*}

\subsection{$n=2$}
For $p=0$, the absolute cutoff is $W=3$:
\[\alpha_{2}^t=\sum_{d=3}^{3}R_{2,d}+O(\mathrm{SR}^{4}).\]
\paragraph{Absolute order 3.}
\begin{align*}
R_{2,3} &=  - \frac{r^{3}}{256 c_{s}^{3}} - \frac{3 \alpha_{0}^{c_s} r^{2}}{64 c_{s}^{2}} \\
&{}  - \frac{3 \alpha_{0}^s r^{2}}{64 c_{s}^{2}} - \frac{\left(\alpha_{0}^{c_s}\right)^{2} r}{8 c_{s}} \\
&{}  - \frac{\left(\alpha_{0}^s\right)^{2} r}{8 c_{s}} + \frac{\alpha_{1}^{c_s} r}{8 c_{s}} \\
&{}  + \frac{\alpha_{1}^s r}{8 c_{s}} - \frac{\alpha_{0}^{c_s} \alpha_{0}^s r}{4 c_{s}}
\end{align*}

\subsection{$n=3$}
For $p=0$, the absolute cutoff is $W=4$:
\[\alpha_{3}^t=\sum_{d=4}^{4}R_{3,d}+O(\mathrm{SR}^{5}).\]
\paragraph{Absolute order 4.}
\begin{align*}
R_{3,4} &= \frac{3 r^{4}}{2048 c_{s}^{4}} - \frac{\alpha_{1}^{c_s} r^{2}}{16 c_{s}^{2}} \\
&{}  - \frac{\alpha_{1}^s r^{2}}{16 c_{s}^{2}} + \frac{\alpha_{2}^{c_s} r}{8 c_{s}} \\
&{}  + \frac{\alpha_{2}^s r}{8 c_{s}} + \frac{\left(\alpha_{0}^{c_s}\right)^{3} r}{8 c_{s}} \\
&{}  + \frac{\left(\alpha_{0}^s\right)^{3} r}{8 c_{s}} + \frac{3 \alpha_{0}^{c_s} r^{3}}{128 c_{s}^{3}} \\
&{}  + \frac{3 \alpha_{0}^s r^{3}}{128 c_{s}^{3}} + \frac{7 \left(\alpha_{0}^{c_s}\right)^{2} r^{2}}{64 c_{s}^{2}} \\
&{}  + \frac{7 \left(\alpha_{0}^s\right)^{2} r^{2}}{64 c_{s}^{2}} - \frac{3 \alpha_{0}^{c_s} \alpha_{1}^{c_s} r}{8 c_{s}} \\
&{}  - \frac{3 \alpha_{0}^{c_s} \alpha_{1}^s r}{8 c_{s}} - \frac{3 \alpha_{0}^s \alpha_{1}^{c_s} r}{8 c_{s}} \\
&{}  - \frac{3 \alpha_{0}^s \alpha_{1}^s r}{8 c_{s}} + \frac{3 \alpha_{0}^{c_s} \left(\alpha_{0}^s\right)^{2} r}{8 c_{s}} \\
&{}  + \frac{3 \left(\alpha_{0}^{c_s}\right)^{2} \alpha_{0}^s r}{8 c_{s}} + \frac{7 \alpha_{0}^{c_s} \alpha_{0}^s r^{2}}{32 c_{s}^{2}}
\end{align*}

\end{document}
